\section{A Global Reduced Solver with Posterior Certificates}
\label{sec:certified-solver}

We first isolate the computation principle behind canonical reduction.  This
also clarifies which part is general Hadamard optimization and which part is
specific to the SPD action bundle.

\begin{theorem}[Canonical strongly-convex fiber engine]
\label{thm:canonical-fiber-engine}
Let \(\mathcal C\) be a closed geodesically convex totally geodesic
submanifold of a finite-dimensional Hadamard manifold.  Suppose
\(F:\mathcal C\to\mathbb R\) is \(C^2\), coercive, and
\(\mu\)-strongly geodesically convex.  If
\[
  \operatorname{Hess}F\preceq L_0g
  \quad\text{on}\quad
  \{x:F(x)\le F(x_0)\},
\]
then \(F\) has a unique minimizer \(x_\star\), every complete update segment
of
\begin{equation}
  x_{k+1}
  =\operatorname{Exp}_{x_k}
   \left(-L_0^{-1}\operatorname{grad}F(x_k)\right)
  \label{eq:abstract-fiber-update}
\end{equation}
stays in the current sublevel set, and, with
\(q_{\mu,L}=1-\mu/L_0\),
\begin{align}
  \Delta_k&\le q_{\mu,L}^k\Delta_0,
  &
  d(x_k,x_\star)
  &\le\sqrt{\frac{2\Delta_0}{\mu}}\,q_{\mu,L}^{k/2},
  \label{eq:abstract-fiber-rates}\\
  \norm{\operatorname{grad}F(x_k)}
  &\le\sqrt{2L_0\Delta_0}\,q_{\mu,L}^{k/2}.
  \label{eq:abstract-fiber-gradient-rate}
\end{align}
At every \(x\in\mathcal C\),
\begin{equation}
  d(x,x_\star)\le\frac{\norm{\operatorname{grad}F(x)}}{\mu},
  \qquad
  0\le F(x)-F(x_\star)
  \le\frac{\norm{\operatorname{grad}F(x)}^2}{2\mu}.
  \label{eq:abstract-fiber-certificates}
\end{equation}
The constants in \eqref{eq:abstract-fiber-certificates} are sharp.
\end{theorem}

The first-return sublevel statement is important: it uses a Hessian bound on
the current sublevel itself, rather than first enlarging its radius to cover a
putative update.  For squared distance on a Hadamard manifold with
\(-\kappa\le\operatorname{sec}\le0\), Hessian comparison supplies the sharp
radius-only majorant
\begin{equation}
  L(D)=\psi(\sqrt\kappa D),
  \qquad
  \psi(s)=
  \begin{cases}
    s\coth s,&s>0,\\
    1,&s=0.
  \end{cases}
  \label{eq:hessian-comparison-factor}
\end{equation}
The general descent ingredients are classical; the theorem records the exact
form needed to connect the canonical section to the action-specific result.

\subsection{Explicit SPD residual and nonasymptotic complexity}

For fixed \(B\in\Bact\), retain \(R_0,\mathcal Q_\Sigma,f_B\), and
\(\mathcal R(\Sigma)\) from \cref{sec:action-bundle}, and define
\begin{equation}
  W_B(\Sigma)=\mathcal R_{22}(\Sigma)^\circ.
  \label{eq:solver-residual}
\end{equation}
This is the exact whitened negative gradient:
\begin{equation}
  -\operatorname{grad}f_B(\Sigma)
  =\Sigma^{1/2}W_B(\Sigma)\Sigma^{1/2},
  \qquad
  \norm{\operatorname{grad}f_B(\Sigma)}_\Sigma
  =\fro{W_B(\Sigma)}.
  \label{eq:residual-gradient-identity}
\end{equation}

Given any \(\Sigma_0\in\mathscr P_n\), put
\begin{equation}
  D_0=\dai(R_0,\mathcal Q_{\Sigma_0}),
  \qquad
  L_0=\psi(D_0/\sqrt2),
  \qquad
  q_0=1-L_0^{-1}.
  \label{eq:solver-constants}
\end{equation}
If \(D_0=0\), the initial controller is canonical.  Otherwise define
\begin{equation}
  \chi_0=-\log q_0
  =-\log(1-L_0^{-1})
  \label{eq:solver-discrete-exponent}
\end{equation}
and run
\begin{equation}
  \boxed{
  \Sigma_{k+1}
  =\Sigma_k^{1/2}
   \exp\!\left(L_0^{-1}W_B(\Sigma_k)\right)
   \Sigma_k^{1/2}.}
  \label{eq:certified-solver-update}
\end{equation}

\begin{theorem}[Global action solver and posterior certificates]
\label{thm:certified-action-solver}
Let \(A\) have full column rank, let \(B\in\Bact\), and let
\(\Sigma_\star\) be the hidden coordinate of \(s_A(B)\).  From every
\(\Sigma_0\in\mathscr P_n\), \eqref{eq:certified-solver-update} preserves
determinant one and satisfies
\begin{align}
  f_B(\Sigma_{k+1})
  &\le f_B(\Sigma_k)
  -\frac{1}{2L_0}\fro{W_B(\Sigma_k)}^2,
  \label{eq:solver-descent}\\
  f_B(\Sigma_k)-f_B(\Sigma_\star)
  &\le \frac{D_0^2}{2}q_0^k,
  \label{eq:solver-value-rate}\\
  \dai(\Sigma_k,\Sigma_\star)
  &=\dai(\Phi_A(B,\Sigma_k),s_A(B))
  \le D_0q_0^{k/2},
  \label{eq:solver-distance-rate}\\
  \fro{W_B(\Sigma_k)}
  &\le \sqrt{L_0}\,D_0q_0^{k/2}.
  \label{eq:solver-residual-rate}
\end{align}
The observable residual gives the sharp posterior certificates
\begin{align}
  \dai(\Phi_A(B,\Sigma_k),s_A(B))
  &\le\fro{W_B(\Sigma_k)},
  \label{eq:solver-distance-certificate}\\
  0\le f_B(\Sigma_k)-f_B(\Sigma_\star)
  &\le\frac12\fro{W_B(\Sigma_k)}^2.
  \label{eq:solver-gap-certificate}
\end{align}
For \(D_0>0\), value, controller, and residual tolerances are guaranteed,
respectively, after
\begin{align}
  k&\ge\frac1{\chi_0}
       \pospart{\log\frac{D_0^2}{2\varepsilon_f}},
  \label{eq:solver-value-complexity}\\
  k&\ge\frac2{\chi_0}
       \pospart{\log\frac{D_0}{\varepsilon_x}},
  \label{eq:solver-distance-complexity}\\
  k&\ge\frac2{\chi_0}
       \pospart{\log\frac{\sqrt{L_0}D_0}{\varepsilon_r}}
  \label{eq:solver-residual-complexity}
\end{align}
iterations, with ceilings understood.
\end{theorem}

The current-sublevel curvature comparison uses
\(L_0=\psi(D_0/\sqrt2)\), the sharp dimension-uniform majorant.  No
radius-doubling argument is needed.  The denominators above invert the
discrete factor \(q_0\) exactly; \(1/\chi_0\le L_0\) recovers the simpler
exponential estimates.

\subsection{Active realization and spectral-arithmetic cost}

The operation counts below use a spectral-arithmetic model: dense
factorizations and matrix functions on an \(r\times r\) symmetric matrix cost
\(O(r^3)\), with scalar elementary functions counted at unit cost.  These are
algebraic operation counts, not bit-complexity or floating-point error bounds.

Let
\[
  \mathcal W_A
  =\operatorname{span}(\operatorname{col}A,\operatorname{col}B),
  \quad
  r=\dim\mathcal W_A\le2m,
  \quad
  s=r-m,
  \quad
  q=d-r.
\]
The active-adapted construction in
\cref{eq:active-residual-splitting} represents the hidden state by
\(\Sigma_a\oplus\beta I_q\), with
\(\det\Sigma_a\,\beta^q=1\), and the full residual norm by
\(\fro{W_a}^2+qw_c^2\).  All complement terms are omitted when \(q=0\).
This family is invariant under the exact update.  A general
\(\Sigma_0\in\mathscr P_{d-m}\) need not lie in it, so
\cref{thm:certified-action-solver} gives arbitrary-start convergence while the
operation count below applies to active-compatible starts.  The identity is
the default such start.

\begin{corollary}[Active canonical solver]
\label{cor:active-core-solver}
Let \(\Sigma_{a,0}\in\SPD^s\) and, when present, \(\beta_0>0\) satisfy
\(\det\Sigma_{a,0}\,\beta_0^q=1\); the identity choice
\(\Sigma_{a,0}=I_s\), \(\beta_0=1\) is admissible.  The updates
\begin{equation}
  \Sigma_{a,k+1}
  =\Sigma_{a,k}^{1/2}e^{W_{a,k}/L_0}\Sigma_{a,k}^{1/2},
  \qquad
  \beta_{k+1}=\beta_ke^{w_{c,k}/L_0}
  \label{eq:active-solver-update}
\end{equation}
are exactly the full iteration \eqref{eq:certified-solver-update}.  They
preserve \(P_kA=B\), \(P_k\succ0\), and \(\det P_k=1\), while the stopping
test \(\fro{W_{a,k}}^2+qw_{c,k}^2\le\varepsilon^2\) implies controller
error at most \(\varepsilon\).

Preprocessing costs \(O(dm^2+r^3)\) spectral-arithmetic operations and
\(O(dm+r^2)\) storage.  Each iteration costs \(O(r^3)\).  For \(D_0>0\), the total
spectral-arithmetic
cost is
\begin{equation}
  O\!\left(
  dm^2+
  r^3\left\{
  1+\frac1{\chi_0}
  \pospart{\log\frac{\sqrt{L_0}D_0}{\varepsilon}}
  \right\}\right),
  \qquad r\le2m.
  \label{eq:active-total-complexity}
\end{equation}
The implicit controller application costs \(O(dr+r^2)\).  The same formulas
cover \(q=0\); if \(s=0\), the solver has no variable and the canonical
controller is
\[
  s_A(B)=EME^\top+\rho(I-EE^\top).
\]
\end{corollary}

\subsection{Conditional inexact residuals}

Finite-precision propagation requires certified numerical information that is
separate from the geometric solver theorem.  Assume a matrix-function routine
supplies a rigorous initial radius majorant \(\widehat D_0\ge D_0\), a
symmetric trace-free direction
\(\widetilde W_k=W_B(\Sigma_k)+E_k\), and an error majorant
\(b_k\ge\fro{E_k}\).  Such bounds require a certified matrix-function
implementation; they are not produced by the analysis below
\cite{higham2008functions}.

Fix \(0\le\delta<1\) and require at every accepted step
\begin{equation}
  b_k\le\frac{\delta}{1+\delta}\fro{\widetilde W_k}.
  \label{eq:computable-relative-error-test}
\end{equation}
This observable test implies
\begin{equation}
  \fro{E_k}\le\delta\fro{W_B(\Sigma_k)}.
  \label{eq:relative-residual-error}
\end{equation}
Define
\begin{equation}
  \widehat L_0=\psi(\widehat D_0/\sqrt2),
  \quad
  c_\delta=\left(\frac{1-\delta}{1+\delta}\right)^2,
  \quad
  \widehat q_\delta=1-\frac{c_\delta}{\widehat L_0},
  \quad
  \eta=\frac{1-\delta}{\widehat L_0(1+\delta)^2}.
  \label{eq:inexact-step}
\end{equation}

\begin{corollary}[Certified inexact residual-oracle propagation]
\label{cor:inexact-action-solver}
If the supplied majorants satisfy the preceding conditions, the update
\[
  \Sigma_{k+1}
  =\Sigma_k^{1/2}\exp(\eta\widetilde W_k)\Sigma_k^{1/2}
\]
has nonincreasing objective radius and
\begin{align}
  f_B(\Sigma_k)-f_B(\Sigma_\star)
  &\le\frac{\widehat D_0^2}{2}\widehat q_\delta^k,
  \label{eq:inexact-value-rate}\\
  \dai(\Phi_A(B,\Sigma_k),s_A(B))
  &\le\widehat D_0\widehat q_\delta^{k/2},
  \label{eq:inexact-distance-rate}\\
  \fro{\widetilde W_k}
  &\le\frac{(1+\delta)^2}{1-\delta}
  \sqrt{\widehat L_0}\widehat D_0\widehat q_\delta^{k/2}.
  \label{eq:inexact-residual-rate}
\end{align}
The observed residual also certifies
\begin{equation}
  \dai(\Phi_A(B,\Sigma_k),s_A(B))
  \le\frac{\fro{\widetilde W_k}}{1-\delta},
  \qquad
  f_B(\Sigma_k)-f_B(\Sigma_\star)
  \le\frac{\fro{\widetilde W_k}^2}{2(1-\delta)^2}.
  \label{eq:inexact-certificates}
\end{equation}
\end{corollary}

This corollary certifies propagation of error in the residual direction
\(\widetilde W_k\) when the accepted update is the exact Riemannian exponential
of that direction.  It does not cover roundoff in the exponential update,
factorizations, active-basis construction, determinant maintenance, or stored
iterates, and it is not a bit-complexity theorem for matrix logarithms, roots,
or exponentials.  End-to-end finite-precision certification requires those
additional errors to be enclosed by the implementation.
